## Test 1, Question 3 — Lawn Sprinkler

Nozzle area:
$a = \dfrac{\pi}{4}(0.008)^2$
$= 5.027\times10^{-5}\text{ m}^2$

Flow through each arm:
$Q = aV$
$= 5.027\times10^{-5}\times 10$
$= 5.027\times10^{-4}\text{ m}^3/\text{s}$

Reading the arm radii as $r_A = 0.25\text{ m}$ (arm to nozzle A) and $r_B = 0.20\text{ m}$ (arm to nozzle B), with both jets discharging tangentially in the same rotational sense:

**Torque to hold the arm stationary:**

$T = \rho Q V(r_A+r_B)$
$= 1000(5.027\times10^{-4})(10)(0.25+0.20)$

$T = \boxed{2.26 \text{ N·m}}$

**Speed of free rotation** — steady state requires the net moment of momentum flux (using *relative* velocity $V-\omega r$ at each nozzle) to vanish:

$V(r_A+r_B) = \omega(r_A^2+r_B^2)$

$\omega = \frac{10(0.45)}{0.25^2+0.20^2}$
$= \frac{4.5}{0.1025}$
$= 43.9\text{ rad/s}$

$N = \frac{60\omega}{2\pi}$
$= \boxed{419 \text{ rpm}}$

---

## Test 1, Question 4 — Oil Nozzle Discharge

$D_2 = 100\text{mm}$, $D_3=25\text{mm}$, so:
$V_3 = V_2(D_2/D_3)^2$
$= 16V_2$

Energy equation, tank surface (1) to nozzle exit (3), both at atmospheric pressure, $V_1\approx0$:

$H = \frac{V_3^2}{2g} + 20\frac{V_2^2}{2g}$
$= \frac{V_3^2}{2g}\left[1+\frac{20}{256}\right]$
$= 1.0781\frac{V_3^2}{2g}$

$4.0 = \frac{V_3^2}{2(9.81)}(1.0781)$
$\;\Rightarrow\; V_3 = \boxed{8.53\text{ m/s}}$

$V_2 = V_3/16$
$= 0.533\text{ m/s}$

**Discharge:**

$Q = \frac{\pi}{4}(0.025)^2(8.53)$
$= \boxed{4.19\text{ L/s}}$

**Pressure at base of nozzle** (point 2), applying the energy equation tank → point 2, with the *entire* pipe loss $20V_2^2/2g$ occurring upstream of the nozzle:

$H = \frac{P_2}{\gamma_{oil}} + \frac{V_2^2}{2g} + 20\frac{V_2^2}{2g}$
$= \frac{P_2}{\gamma_{oil}} + 21\frac{V_2^2}{2g}$

$\frac{P_2}{\gamma_{oil}} = 4.0 - 21\frac{(0.533)^2}{2(9.81)}$
$= 4.0-0.305$
$= 3.695\text{ m}$

$P_2 = (0.8\times9810)(3.695)$
$= \boxed{29.0\text{ kPa (gauge)}}$

---

## Test 2, Question 1(a)

Continuity (incompressible):
$\dfrac{\partial u}{\partial x}+\dfrac{\partial v}{\partial y}=0$

$\frac{\partial u}{\partial x} = 2Ax$
$\;\Rightarrow\; \frac{\partial v}{\partial y} = -2Ax$

$v = -2Axy + f(x)$

Applying $v=0$ at $y=0$ gives $f(x)=0$:

$\boxed{v = -2Axy}$

**Irrotationality** — check $\omega_z = \dfrac{\partial v}{\partial x}-\dfrac{\partial u}{\partial y}$:

$\frac{\partial v}{\partial x} = -2Ay$

$\frac{\partial u}{\partial y} = B$

$\omega_z = -2Ay - B \neq 0 \text{ in general}$

**The flow is rotational** (not irrotational) unless $A=B=0$.

---

## Test 2, Question 1(b)

$u=x^2y$
$v=y^2z$
$w=-(2xyz+yz^2)$

**Continuity check:**

$\frac{\partial u}{\partial x}+\frac{\partial v}{\partial y}+\frac{\partial w}{\partial z}$
$= 2xy + 2yz -(2xy+2yz)$
$= 0$

Continuity is satisfied everywhere, and the field has no explicit time dependence → **this is a possible case of steady, incompressible flow.**

**Velocity at (2,1,3):**

$u=4$
$v=3$
$w=-(12+9)=-21$

$\vec V = 4\hat i+3\hat j-21\hat k \text{ m/s}$

$\vert{}\vec V\vert{} = \boxed{21.6\text{ m/s}}$

**Acceleration** (steady flow → purely convective):

$a_x = u\frac{\partial u}{\partial x}+v\frac{\partial u}{\partial y}+w\frac{\partial u}{\partial z}$
$= 2x^3y^2+x^2y^2z$
$= 16+12$
$= 28$

$a_y = u\frac{\partial v}{\partial x}+v\frac{\partial v}{\partial y}+w\frac{\partial v}{\partial z}$
$= 2y^3z^2+wy^2$
$= 18-21$
$= -3$

$a_z = u\frac{\partial w}{\partial x}+v\frac{\partial w}{\partial y}+w\frac{\partial w}{\partial z}$
$= -24-63+210$
$= 123$

$\vec a = 28\hat i-3\hat j+123\hat k\text{ m/s}^2$

$\vert{}\vec a\vert{} = \boxed{126.2\text{ m/s}^2}$

---

## Test 2, Question 2(a) — Pipe Bend

$A = \dfrac{\pi}{4}(0.3)^2$
$= 0.07069\text{ m}^2$

$Q = AV$
$= 0.2474\text{ m}^3/\text{s}$

Since diameter (hence $V$) is unchanged through the bend and no loss is stated:

$P_1=P_2 = \rho g H$
$= 1000(9.81)(20)$
$= 196.2\text{ kPa}$

**x-direction:**

$F_x = (1-\cos45^\circ)(PA+\rho QV)$
$= 0.2929(13{,}869+866)$
$= \boxed{4315\text{ N}}$

**y-direction:**

$F_y = \sin45^\circ(PA-\rho QV)$
$= 0.7071(13{,}869-866)$
$= \boxed{9196\text{ N}}$

**Resultant:**

$R = \sqrt{F_x^2+F_y^2}$
$= \boxed{10.16\text{ kN}}$

$\theta = \tan^{-1}\left(\frac{F_y}{F_x}\right)$
$= \boxed{64.9^\circ \text{ from the initial pipe axis}}$

---

## Test 2, Question 2(b) — Conical Diffuser

Diameter varies linearly:
$D(x) = 0.20+0.30x\text{ (m)}$

At $x=1.5\text{ m}$:
$D=0.65\text{ m}$

$A=\dfrac{\pi}{4}(0.65)^2$
$=0.3318\text{ m}^2$

At the instant:
$Q=0.2\text{ m}^3/\text{s}$

$\dfrac{dQ}{dt}=0.05\text{ m}^3/\text{s}^2$

$V = Q/A$
$= 0.6027\text{ m/s}$

**Local acceleration:**

$\frac{\partial V}{\partial t} = \frac{1}{A}\frac{dQ}{dt}$
$= \frac{0.05}{0.3318}$
$= \boxed{0.151\text{ m/s}^2}$

**Convective acceleration** — with $\dfrac{dA}{dx} = \dfrac{\pi}{2}D\dfrac{dD}{dx} = \dfrac{\pi}{2}(0.65)(0.30)=0.3063$:

$\frac{\partial V}{\partial x} = -\frac{Q}{A^2}\frac{dA}{dx}$
$= -\frac{0.2}{0.1101}(0.3063)$
$= -0.5565\text{ s}^{-1}$

$V\frac{\partial V}{\partial x} = 0.6027(-0.5565)$
$= \boxed{-0.336\text{ m/s}^2}$

**Total acceleration:**

$a = \frac{\partial V}{\partial t}+V\frac{\partial V}{\partial x}$
$= 0.151-0.336$
$= \boxed{-0.185\text{ m/s}^2}$
